\documentclass{article}


% if you need to pass options to natbib, use, e.g.:
%     \PassOptionsToPackage{numbers, compress}{natbib}
% before loading neurips_2024


% ready for submission
\usepackage{neurips_2025}
\usepackage{amsmath}

% to compile a preprint version, e.g., for submission to arXiv, add add the
% [preprint] option:
%     \usepackage[preprint]{neurips_2024}


% to compile a camera-ready version, add the [final] option, e.g.:
%     \usepackage[final]{neurips_2024}


% to avoid loading the natbib package, add option nonatbib:
%    \usepackage[nonatbib]{neurips_2024}


\usepackage[utf8]{inputenc} % allow utf-8 input
\usepackage[T1]{fontenc}    % use 8-bit T1 fonts
\usepackage{hyperref}       % hyperlinks
\usepackage{url}            % simple URL typesetting
\usepackage{booktabs}       % professional-quality tables
\usepackage{amsfonts}       % blackboard math symbols
\usepackage{nicefrac}       % compact symbols for 1/2, etc.
\usepackage{microtype}      % microtypography
\usepackage{xcolor}         % colors



\usepackage{algorithm}
\usepackage{algorithmic}
\usepackage{wrapfig}
\usepackage{graphicx}
\usepackage{enumerate}
\usepackage{multirow}
\usepackage{wrapfig}
\usepackage{lipsum}
\usepackage{xcolor}
\usepackage{diagbox} 
\usepackage{enumitem}
\usepackage{subcaption}


\usepackage{amsmath,amssymb,mathtools}

\newcommand{\R}{\mathbb{R}}
\newcommand{\E}{\mathbb{E}}
\newcommand{\norm}[1]{\left\lVert #1 \right\rVert}
\newcommand{\inner}[2]{\left\langle #1, #2 \right\rangle}
\newcommand{\sigmoid}{\sigma}
\newcommand{\tr}{\mathrm{tr}}
\newcommand{\eps}{\varepsilon}


\title{Manifold-of-Thought: A Training-Free State-Driven Reasoning Control Paradigm for Large Language Models}


% The \author macro works with any number of authors. There are two commands
% used to separate the names and addresses of multiple authors: \And and \AND.
%
% Using \And between authors leaves it to LaTeX to determine where to break the
% lines. Using \AND forces a line break at that point. So, if LaTeX puts 3 of 4
% authors names on the first line, and the last on the second line, try using
% \AND instead of \And before the third author name.


\author{%
  David S.~Hippocampus\thanks{Use footnote for providing further information
    about author (webpage, alternative address)---\emph{not} for acknowledging
    funding agencies.} \\
  Department of Computer Science\\
  Cranberry-Lemon University\\
  Pittsburgh, PA 15213 \\
  \texttt{hippo@cs.cranberry-lemon.edu} \\
  % examples of more authors
  % \And
  % Coauthor \\
  % Affiliation \\
  % Address \\
  % \texttt{email} \\
  % \AND
  % Coauthor \\
  % Affiliation \\
  % Address \\
  % \texttt{email} \\
  % \And
  % Coauthor \\
  % Affiliation \\
  % Address \\
  % \texttt{email} \\
  % \And
  % Coauthor \\
  % Affiliation \\
  % Address \\
  % \texttt{email} \\
}


\begin{document}


\maketitle

\begin{abstract}
Large Language Models (LLMs) benefit from intermediate reasoning, but most existing methods still impose a fixed external script (e.g., linear CoT, sampling, or tree search) before inference starts. This makes the reasoning procedure rigid and weakly coupled to the model's evolving internal logic.

We introduce \textbf{Manifold-of-Thought (MoT)}, a \emph{training-free} state-driven reasoning control framework. MoT does not train a policy network, does not rely on dense head-level control, and does not treat internal state as a correctness classifier. Instead, MoT maintains a compact sentence-level logic state
$m_t=[\delta_t, v_t, c_t, H_t, e_t]$ that captures local structure, stepwise progress, directional consistency, uncertainty, and a slow-time stagnation debt. From mixed-source trajectories, MoT builds soft local reasoning regimes and estimates regime-conditioned action preferences. At each step, it selects structured actions through
\[
\text{state} \rightarrow \text{soft regime} \rightarrow \text{node action} \rightarrow \text{memory routing},
\]
with hard structural constraints and without policy learning. Here, \emph{training-free} means no gradient-based optimization of either the backbone LLM or an auxiliary controller during adaptation; offline estimation is limited to unsupervised regime fitting and nonparametric preference statistics.

The action space is intentionally small: \texttt{continue}, \texttt{stop}, \texttt{revisit}, \texttt{summarize}, and \texttt{branch}; memory routing uses only \texttt{recent}, \texttt{anchor+recent}, \texttt{summary+recent}, and \texttt{branch-local}. This yields a simple, paradigm-level closed loop that supports both serial reasoning and budgeted branching while staying practical for long-horizon and revision-heavy tasks.
\end{abstract}


\section{Introduction}
\label{sec:intro}

Reasoning support is now central for LLM performance, but most methods still hard-code an external procedure in advance: linear chains, repeated sampling, or explicit tree search. These approaches can be effective, yet they keep the control logic outside the model's evolving internal state.

We focus on a different question: can reasoning be driven by \emph{local internal state} rather than fixed scripts or learned controllers?

Our answer is \textbf{Manifold-of-Thought (MoT)} as a state-driven, training-free \emph{reasoning control} framework. MoT does not introduce another fixed prompting script. Instead, at sentence boundaries it reads compact hidden-state dynamics, infers the current local reasoning regime, and triggers a small set of structured actions on node growth and memory routing.

The key framing is deliberately weak and practical: internal state should \emph{not} be expected to reliably predict final correctness at each step. It only needs to identify the current local regime well enough to prefer statistically better next actions under that regime.

This leads to a closed loop:
\[
\text{hidden-state trajectory}
\rightarrow
m_t
\rightarrow
R_t
\rightarrow
(a_t^{\text{node}}, a_t^{\text{mem}}),
\]
where $m_t$ is a compact logic state, $R_t$ is a soft regime assignment, and actions are executed by an external actuator.

Compared with controller-style methods, MoT makes three simplifying commitments: (i) no policy learning, (ii) no dense head-control as the primary interface, and (iii) no heavy semantic role tagging for memory. The result is a compact, deployable control paradigm that preserves adaptivity while reducing unnecessary complexity.

\paragraph{Contributions.}
\begin{itemize}[leftmargin=1.3em, itemsep=0.2em]
    \item \textbf{A paradigm shift from script design to state-driven control.} We formulate reasoning as a sentence-level state-action loop driven by internal logic trajectory rather than by fixed external scripts.
    \item \textbf{A minimal and robust logic state.} We use a 5D compact state $m_t=[\delta_t,v_t,c_t,H_t,e_t]$ where $e_t$ is a slow-time debt variable that captures multi-step stagnation and drift not visible from one-step features.
    \item \textbf{Training-free soft-regime action selection.} We infer soft regimes from trajectory coverage and choose structured actions using regime preferences plus structural constraints, without RL or supervised policy training.
    \item \textbf{A practical node/memory actuator design.} MoT unifies node-level operations (\texttt{continue}/\texttt{stop}/\texttt{revisit}/\texttt{summarize}/\texttt{branch}) with sparse memory routing templates, enabling adaptive yet controlled long-horizon reasoning.
\end{itemize}


\section{Preliminaries}
\label{sec:prelim}

\subsection{Sentence-Level Reasoning Process}
Consider an autoregressive transformer LLM with $L$ layers, $H$ attention heads per layer, and hidden size $D$. 
For token $i$ at layer $\ell$, denote hidden state by
\[
h^{(\ell)}_{i} \in \R^{D}.
\]

We decompose reasoning into sentence-level units
\[
y_{1}, y_{2}, \dots, y_{T},
\]
where each $y_t$ is generated between two sentence boundaries under a fixed delimiter policy shared by offline and online stages. In our minimal instantiation, a step ends at newline or sentence-ending punctuation (\texttt{.?!}); if no delimiter appears, we force a boundary at $N_{\max}^{\mathrm{step}}$ tokens.

At step $t$, MoT extracts a compact state $m_t$, infers soft regimes, selects node and memory actions, and executes them before generating $y_{t+1}$:
\[
y_t \rightarrow m_t \rightarrow R_t \rightarrow (a_t^{\text{node}}, a_t^{\text{mem}}) \rightarrow y_{t+1}.
\]

This sentence-level granularity balances responsiveness and stability, and keeps the action interface interpretable and implementable.

\subsection{Core Assumption (Weak Form)}
We do not assume that internal states can reliably classify final correctness at each step. We only assume that LLM reasoning repeatedly visits local regimes whose state signatures are statistically associated with more suitable next actions. MoT therefore performs \emph{regime-aware action selection}, not correctness prediction.

\subsection{Training-Free Scope}
In this paper, \emph{training-free} means: no gradient updates to the backbone LLM, and no trained policy/controller during adaptation. Allowed offline operations are unsupervised regime fitting (e.g., GMM/prototype kernels), nonparametric preference estimation, and validation-time hyperparameter selection.


\section{Manifold-of-Thought (MoT)}
\label{sec:method}

\subsection{Overview}
MoT is a training-free state-driven control loop over structured reasoning actions. At each step, the system:
\begin{enumerate}[leftmargin=1.3em, itemsep=0.2em]
    \item extracts sentence-level state features from hidden-state dynamics,
    \item updates compact logic state $m_t$,
    \item computes soft regime assignment $\pi_k(m_t)=p(R_k|m_t)$,
    \item selects node action $a_t^{\text{node}}$ with structural constraints,
    \item derives memory routing mode $a_t^{\text{mem}}$ and updates active context.
\end{enumerate}

No policy network is trained. No dense head-control module is required by the core method.

\paragraph{Core principle vs. instantiation.}
The \textbf{core MoT principle} has four components: state extractor, regime estimator, action scorer, and actuator.
This paper provides one \textbf{minimal instantiation}: 5D state ($[\delta,v,c,H,e]$), five node actions, four memory templates, and budgeted structural constraints.

\subsection{Sentence-Level State Extraction}
\label{sec:obs}
At the end of sentence $y_t$, let $n_t$ be its number of tokens. 
For token $i$ in $y_t$, we use top-layer hidden state:
\begin{equation}
h_{t,i} \;=\; h^{(L)}_{t,i},
\qquad h_{t,i}\in\R^D,
\label{eq:lastk}
\end{equation}
with optional last-$k$ average as an ablation.

Define sentence center
\begin{equation}
z_t \;=\; \frac{1}{n_t}\sum_{i=1}^{n_t} h_{t,i} \in \R^D,
\label{eq:zt}
\end{equation}
within-sentence dispersion
\begin{equation}
\delta_t \;=\; \tr\!\left(
\frac{1}{n_t}\sum_{i=1}^{n_t}(h_{t,i}-z_t)(h_{t,i}-z_t)^\top
\right)
\;=\; \frac{1}{n_t}\sum_{i=1}^{n_t} \norm{h_{t,i}}_2^2 - \norm{z_t}_2^2.
\label{eq:disp}
\end{equation}
progress magnitude
\begin{equation}
v_t \;=\; \norm{z_t-z_{t-1}}_2,
\label{eq:speed}
\end{equation}
directional consistency
\begin{equation}
c_t \;=\;
\begin{cases}
1, & t=1,\\[3pt]
\dfrac{\inner{z_t-z_{t-1}}{z_{t-1}-z_{t-2}}}
{\norm{z_t-z_{t-1}}_2\norm{z_{t-1}-z_{t-2}}_2+\eps}, & t\ge 2,
\end{cases}
\label{eq:consistency}
\end{equation}
and sentence uncertainty
\begin{equation}
H_t \;=\; \frac{1}{n_t}\sum_{i=1}^{n_t} \mathcal{H}(p_{t,i}),
\label{eq:entropy}
\end{equation}
where $p_{t,i}$ is the next-token distribution at token $i$.

For fixed-scale robustness, we use source statistics $(\mu,\sigma)\in\R^4$ to normalize
\begin{equation}
\hat s_t = \frac{[\delta_t,v_t,c_t,H_t]^\top-\mu}{\sigma+\eps}.
\label{eq:norm}
\end{equation}

\subsection{Slow-Time Debt Variable and Final Logic State}
\label{sec:mt}
Single-step features cannot capture persistent stagnation or drift. We introduce a slow variable
\begin{equation}
e_t = \alpha e_{t-1} + (1-\alpha)u_t,
\label{eq:debt}
\end{equation}
where local debt $u_t$ is
\begin{equation}
u_t =
\lambda_v\,\mathrm{lowprog}(v_t)+
\lambda_c\,\mathrm{turn}(c_t)+
\lambda_H\,H_t.
\label{eq:ut}
\end{equation}
We use simple clipped forms:
\begin{equation}
\mathrm{lowprog}(v_t)=\max\!\left(0,\frac{\tau_v-v_t}{\tau_v+\eps}\right),\qquad
\mathrm{turn}(c_t)=\max(0,-c_t),
\label{eq:debt_terms}
\end{equation}
so debt rises under under-progress and directional reversal.

The compact logic state is
\begin{equation}
m_t=[\hat s_t^\top,e_t]^\top \in \R^5.
\label{eq:mt}
\end{equation}
To keep the fifth dimension comparable with $\hat s_t$, we use bounded debt normalization:
\begin{equation}
\bar e_t=\mathrm{clip}(e_t,0,e_{\max})/e_{\max},\qquad
m_t=[\hat s_t^\top,\bar e_t]^\top.
\label{eq:enorm}
\end{equation}

\subsection{Soft Regimes from State Coverage}
\label{sec:regime}
Regimes are latent local reasoning situations, not hand-labeled script names. From mixed-source trajectories
\[
\mathcal{D}_{\mathrm{mix}}=\bigcup_{j=1}^{J}\mathcal{D}_j,
\]
we collect sentence-level states and fit a soft partition (e.g., GMM or prototype kernels):
\begin{equation}
\pi_k(m_t)=p(R_k\mid m_t),\qquad k=1,\dots,K.
\label{eq:softregime}
\end{equation}
In practice we use a compact $K=5$ regime set (exploratory/progressing/redirecting/stagnant/converging) only as interpretation names; the model uses soft weights $\pi_k$ directly.

\subsection{Action Space and Memory Routing}
\label{sec:action}
MoT uses two coupled decisions:
\[
a_t^{\text{node}} \in \{\texttt{continue},\texttt{stop},\texttt{revisit},\texttt{summarize},\texttt{branch}\},
\]
\[
a_t^{\text{mem}} \in \{\texttt{recent},\texttt{anchor+recent},\texttt{summary+recent},\texttt{branch-local}\}.
\]

\paragraph{Node actions.}
\texttt{continue}: extend current frontier; 
\texttt{stop}: enter answer-commit mode (one final synthesis step, no further structural expansion);
\texttt{revisit}: rollback to an anchor and regrow suffix;
\texttt{summarize}: compress local chain or merge branches into summary node;
\texttt{branch}: open a sibling continuation with strict budget/depth limits.

\paragraph{Memory routing templates.}
\texttt{recent}: latest window only;
\texttt{anchor+recent}: anchor node plus latest window;
\texttt{summary+recent}: summary node plus latest window;
\texttt{branch-local}: current branch chain plus parent anchor.

Actions are layered, not a single joint mega-label:
\[
m_t \rightarrow R_t \rightarrow a_t^{\text{node}} \rightarrow a_t^{\text{mem}}.
\]
Default mapping is
\texttt{continue}$\rightarrow$\texttt{recent},
\texttt{revisit}$\rightarrow$\texttt{anchor+recent},
\texttt{summarize}$\rightarrow$\texttt{summary+recent},
\texttt{branch}$\rightarrow$\texttt{branch-local}.

\subsection{Node Graph and Active Context}
\label{sec:graph}
Each sentence corresponds to a node with fields
\[
(\text{id},\text{content},\text{parent},\text{depth},\text{step},\text{branch},\text{is\_summary},\text{is\_anchor},\text{is\_terminal}).
\]
At step $t$, only an active context set is fed to decoding:
\[
M_t \subseteq \{y_1,\dots,y_t\},
\]
selected by $a_t^{\text{mem}}$. In this paper, the minimal implementation uses \emph{explicit context assembly} at sentence boundaries. KV selection and sparse masking are treated as future system-level implementations. Each memory template uses a fixed token budget $B_{\mathrm{ctx}}$: include required anchor/summary nodes first, then fill the remaining budget with newest nodes in reverse chronological order.

\paragraph{Anchor policy (minimal).}
To make \texttt{revisit} executable, we define valid anchors by fixed rules: prompt root is always an anchor; every summary node is an anchor; and stable low-debt nodes ($\bar e_t<\tau_{\mathrm{anchor}}$, $c_t>\tau_{c,\mathrm{anchor}}$) are anchor candidates. By default, revisit selects the most recent valid anchor unless an earlier anchor has higher structural score. Ties are broken deterministically by higher depth, then more recent step, then smaller node id.

\paragraph{Actuator contract consistency.}
All actions are defined relative to a fixed actuator contract. The same action semantics, node-transition rules, and memory-routing implementation are required in offline log derivation and online execution.

\subsection{Offline Preference Construction}
\label{sec:offline}
Mixed external reasoning modes are used only for state coverage, not as direct action labels. Action labels are derived from structural transitions and actuator logs:
\texttt{continue}/\texttt{stop}/\texttt{revisit}/\texttt{summarize}/\texttt{branch}.

We separate \textbf{regime coverage} from \textbf{action support}: broad mixed-source traces cover state space, while action preferences are estimated only from actuator-compatible logs. Actions with insufficient support in a regime are smoothed or disabled for that regime.

We keep both successful and unsuccessful trajectories for regime coverage, and use quality-weighted action statistics. Step weight is
\begin{equation}
w_t=\alpha_q\,\tilde q(x)+\alpha_p\,\tilde p_t(x)-\alpha_c\,\tilde c_t(x),\qquad w_t\leftarrow \max(0,w_t),
\label{eq:wt}
\end{equation}
where $\tilde q,\tilde p_t,\tilde c_t$ are normalized outcome quality, prefix utility, and cost terms.
In the minimal implementation:
\begin{itemize}[leftmargin=1.3em, itemsep=0.2em]
    \item $\tilde q(x)\in[0,1]$: final outcome score (e.g., exact match / task success),
    \item $\tilde p_t(x)\in[-1,1]$: short-horizon improvement proxy after step $t$,
    \item $\tilde c_t(x)\in[0,1]$: normalized cumulative compute cost up to step $t$.
\end{itemize}
We instantiate $\tilde p_t$ by a short-window state-improvement proxy:
\begin{equation}
\tilde p_t = \beta_1\,\Delta(-e)_{t\rightarrow t+h} + \beta_2\,\Delta c_{t\rightarrow t+h} + \beta_3\,\Delta(-H)_{t\rightarrow t+h},
\label{eq:pt}
\end{equation}
where $h$ is a small horizon (default $h=1$ sentence).

Because logs are observational, we estimate a \emph{regime-conditioned preference proxy} (not unbiased causal value):
\begin{equation}
U_k(a) =
\frac{\sum_{(x,t)} w_t\,\pi_k(m_t)\,\mathbf{1}[a_t=a]}
{\sum_{(x,t)} \pi_k(m_t)\,\mathbf{1}[a_t=a]+\eps}.
\label{eq:uk}
\end{equation}
Support is
\begin{equation}
N_k(a)=\sum_{(x,t)} \pi_k(m_t)\,\mathbf{1}[a_t=a].
\label{eq:support}
\end{equation}
If $N_k(a)<n_{\min}$, we apply shrinkage to a prior or mark $a$ infeasible under regime $k$. This preserves bad-state distribution while preventing low-quality traces from dominating preferences.
\begin{equation}
\hat U_k(a)=\frac{N_k(a)}{N_k(a)+\lambda_s}U_k(a)+\frac{\lambda_s}{N_k(a)+\lambda_s}U_0(a),
\label{eq:shrink}
\end{equation}
where $U_0(a)$ is a global action prior estimated from all logs. We use $\hat U_k(a)$ for scoring when $N_k(a)<n_{\min}$, and disable action $a$ in regime $k$ when $N_k(a)<n_{\mathrm{hard}}$.
As a robustness variant, one can replace Eq.~\eqref{eq:pt} with final-outcome-only weighting and report sensitivity in ablations.

\subsection{Training-Free Online Decision}
\label{sec:decision}
MoT does not learn $m_t\!\rightarrow\!a_t$ as a policy network. Online action is selected by
\begin{equation}
\mathrm{Score}(a) = \sum_{k=1}^{K}\pi_k(m_t)\,\tilde U_k(a) + B_{\mathrm{struct}}(a,G_t),
\label{eq:score}
\end{equation}
\begin{equation}
a_t^{\text{node}}=\arg\max_a \mathrm{Score}(a),
\label{eq:argmax}
\end{equation}
where $\tilde U_k(a)=\hat U_k(a)$ if $N_k(a)<n_{\min}$ and $\tilde U_k(a)=U_k(a)$ otherwise. $G_t$ is the current node graph and $B_{\mathrm{struct}}$ enforces structural constraints (hard masks or penalties), e.g., no \texttt{revisit} without valid anchor, no \texttt{branch} after budget exhaustion, no \texttt{summarize} below minimal span.
For implementation, we use a feasible action set $\mathcal{A}_{\mathrm{feas}}(G_t)$ and set $\mathrm{Score}(a)=-\infty$ for $a\notin\mathcal{A}_{\mathrm{feas}}(G_t)$.

\subsection{Adaptive Stop and Structural Constraints}
\label{sec:termination}
Stopping is integrated into the same action space. We permit \texttt{stop} only after $t\ge t_{\min}$ and under one of:
\begin{equation}
H_t<\tau_H,\quad v_t<\tau_v,\quad c_t>\tau_c
\label{eq:stop_converge}
\end{equation}
(converging state), or
\begin{equation}
\bar e_t>\tau_e \ \text{and no higher-scoring feasible non-stop action}.
\label{eq:stop_stagnant}
\end{equation}
Here \texttt{stop} means ``switch to answer-commit mode'' rather than ``state certifies correctness''.
To reduce oscillation and early-stop noise, we use hysteresis: stop-entry thresholds are stricter than stop-exit thresholds, and Eq.~\eqref{eq:stop_converge} must hold for $r_{\mathrm{stop}}$ consecutive steps.

To avoid loops:
\begin{itemize}[leftmargin=1.3em, itemsep=0.2em]
    \item per-sample \texttt{revisit} budget,
    \item each anchor revisited at most once,
    \item if post-revisit returns quickly to same regime, disable repeated revisit.
\end{itemize}
\texttt{Branch} is disabled by default and only enabled when all hold: exploratory regime mass above threshold, remaining branch budget/depth, and sufficient support $N_k(\texttt{branch})\ge n_{\min}$.

\subsection{MoT Inference Algorithm}
\label{sec:algo}

Algorithm~\ref{alg:mot} summarizes the full online loop. Offline components are state normalization, soft regimes, and $U_k(a)$ statistics.

\begin{algorithm}[t]
\caption{MoT Inference (State-Driven Node/Memory Control)}
\label{alg:mot}
\begin{algorithmic}[1]
\STATE \textbf{Input:} prompt $x$, frozen LLM
\STATE \textbf{Precomputed:} normalization stats $(\mu,\sigma)$, soft regime model $\pi_k(\cdot)$, utilities $U_k(a)$ and support counts $N_k(a)$
\STATE Initialize node graph $G_0$ with prompt root; initialize active context $M_0$
\STATE Initialize $(z_0,\delta_0,v_0,c_0,H_0)$ from prompt; set $e_0=0$, $\bar e_0=0$, and $m_0=[\hat s_0^\top,\bar e_0]^\top$
\FOR{$t=0,1,\dots,T_{\max}-1$}
    \STATE Compute soft regimes $\pi_k(m_t)$ by Eq.~\eqref{eq:softregime}
    \STATE Compute support-aware feasible action scores by Eq.~\eqref{eq:score}; choose $a_t^{\text{node}}$ by Eq.~\eqref{eq:argmax}
    \STATE Derive $a_t^{\text{mem}}$ from node action and constraints
    \IF{$a_t^{\text{node}}=\texttt{stop}$}
        \STATE Run answer-commit mode (single final synthesis step) and output answer; \textbf{break}
    \ENDIF
    \STATE Execute actuator on graph $G_t$ for non-stop action (continue / revisit / summarize / branch)
    \STATE Update active context to $M_{t+1}$ according to $a_t^{\text{mem}}$
    \STATE Generate next reasoning sentence $y_{t+1}$ under updated context
    \STATE Extract $(\delta_{t+1},v_{t+1},c_{t+1},H_{t+1})$ and update $e_{t+1}$ by Eq.~\eqref{eq:debt}
    \STATE Normalize to $\hat s_{t+1}$ by Eq.~\eqref{eq:norm}; form $m_{t+1}$ by Eq.~\eqref{eq:enorm}
\ENDFOR
\STATE \textbf{if} no explicit stop by $T_{\max}$: force final answer decoding
\end{algorithmic}
\end{algorithm}

\subsection{Minimal Implementation Contract}
\label{sec:contract}
To ensure one-to-one mapping from method to code, a MoT implementation should expose the following minimal interfaces:
\begin{itemize}[leftmargin=1.3em, itemsep=0.2em]
    \item \textbf{State extractor:} given generated sentence $y_t$, return $(\delta_t,v_t,c_t,H_t,\hat s_t,e_t,m_t)$.
    \item \textbf{Regime estimator:} given $m_t$, return soft weights $\{\pi_k(m_t)\}_{k=1}^K$.
    \item \textbf{Action scorer:} given $(m_t,G_t)$, return scores for feasible node actions using Eq.~\eqref{eq:score}.
    \item \textbf{Actuator:} apply \texttt{continue}/\texttt{revisit}/\texttt{summarize}/\texttt{branch}/\texttt{stop} to the node graph.
    \item \textbf{Memory router:} construct $M_t$ from $a_t^{\text{mem}}$ using one of the four templates.
    \item \textbf{Trajectory logger:} log per-step $(m_t,\pi_t,a_t^{\text{node}},a_t^{\text{mem}},w_t,\text{cost})$ for offline preference updates and ablations.
\end{itemize}
The logger must use the same actuator contract as online inference.
The reference implementation must use explicit context assembly (not KV-level routing) for reproducibility.

The required hyperparameter set is intentionally small:
\[
\{K,\alpha,\lambda_v,\lambda_c,\lambda_H,\tau_v,\tau_H,\tau_c,\tau_e,e_{\max},t_{\min},r_{\mathrm{stop}},B_{\mathrm{ctx}},B_{\mathrm{branch}},D_{\mathrm{branch}},B_{\mathrm{revisit}},n_{\min},n_{\mathrm{hard}},\lambda_s\},
\]
where $B_{\mathrm{ctx}}$ is context token budget, $B_{\mathrm{branch}}$ is branch budget, $D_{\mathrm{branch}}$ is maximum branch depth, and $B_{\mathrm{revisit}}$ is revisit budget. All are calibrated on source validation and then frozen for transfer evaluation.

\section{Experiments}
\label{sec:experiments}

\subsection{Evaluation Goals}
Experiments test four claims:
\begin{enumerate}[leftmargin=1.3em, itemsep=0.2em]
    \item \textbf{State-driven adaptivity:} regime-aware actions improve reasoning dynamics over fixed scripts.
    \item \textbf{Efficiency:} state-driven stop/summarize/memory routing reduces cost.
    \item \textbf{Quality:} final accuracy remains competitive or improves.
    \item \textbf{Robustness:} structural constraints prevent loop-like failure modes.
\end{enumerate}

\subsection{Tasks and Benchmarks}
We evaluate on benchmark families emphasizing complementary stressors:
\begin{itemize}[leftmargin=1.3em, itemsep=0.2em]
    \item \textbf{Sequential derivation} (multi-step structured reasoning),
    \item \textbf{Revision-sensitive tasks} (need backtracking and redirection),
    \item \textbf{Long-horizon tasks} (prone to overthinking, drift, and context overload).
\end{itemize}
We instantiate each family with standard benchmarks and keep model, prompting, and evaluation protocols consistent across methods.

\subsection{Experimental Setup}
The backbone LLM is frozen throughout. We collect mixed-source trajectories, compute normalized states, fit soft regimes, and estimate observational preference proxies with quality weighting and support-aware shrinkage (Eq.~\eqref{eq:wt}--\eqref{eq:shrink}). No RL, no supervised policy, and no gradient updates are used.
Transfer protocol is fixed: state extractor and regime estimator are frozen across datasets; only optional temperature scaling on action scores is allowed on a small validation split.

Action constraints (branch/revisit budgets, summarize minimal span, valid-anchor checks) are fixed from source validation and then frozen across target benchmarks. This keeps the method training-free and reproducible.

\subsection{Baselines}
\begin{itemize}[leftmargin=1.3em, itemsep=0.2em]
    \item \textbf{Vanilla decoding:} direct answer generation.
    \item \textbf{Standard CoT:} fixed linear reasoning.
    \item \textbf{Budgeted ToT-lite:} limited branching search baseline.
    \item \textbf{MoT w/o regime softness:} hard nearest-regime action lookup.
    \item \textbf{MoT w/o debt $e_t$:} remove slow-time variable.
    \item \textbf{MoT w/o memory routing:} node actions only.
    \item \textbf{MoT w/o structural constraints:} unconstrained action selection.
    \item \textbf{MoT full:} state-driven soft regime + node/memory actions + constraints.
\end{itemize}

\subsection{Metrics and Pareto Frontiers}
We report quality-cost trade-offs:
\begin{itemize}[leftmargin=1.3em, itemsep=0.2em]
    \item \textbf{Quality:} accuracy / exact match / task success.
    \item \textbf{Cost:} total tokens, number of reasoning nodes, wall-clock latency.
    \item \textbf{Control statistics:} revisit rate, summarize rate, branch usage, average active-context size.
\end{itemize}
In addition to Pareto frontiers, we include matched-budget comparisons.

\subsection{Ablations}
Key ablations:
\begin{itemize}[leftmargin=1.3em, itemsep=0.2em]
    \item \textbf{State dimension:} compare 4D ($[\delta,v,c,H]$) vs.\ 5D ($[\delta,v,c,H,e]$).
    \item \textbf{Regime modeling:} GMM vs.\ prototype-kernel vs.\ hard clustering.
    \item \textbf{Action policy form:} weighted score (Eq.~\eqref{eq:score}) vs.\ rule-only.
    \item \textbf{Constraint strength:} full constraints vs.\ relaxed budgets.
    \item \textbf{Memory templates:} full set vs.\ \texttt{recent}-only.
    \item \textbf{Source mixture composition:} CoT-only vs.\ mixed trajectories for regime coverage.
\end{itemize}

\subsection{Interpretability}
We complement the main experiments with interpretability analyses:
\begin{enumerate}[leftmargin=1.3em, itemsep=0.2em]
    \item \textbf{Logic trajectories:} visualize $m_t$ paths and regime transitions.
    \item \textbf{Event overlays:} highlight high uncertainty, turning, and high-debt segments.
    \item \textbf{Action traces:} inspect where revisit/summarize/branch are triggered.
    \item \textbf{Context traces:} show active-context evolution under memory routing.
\end{enumerate}


\section{Discussion and Potential Limitations}
\label{sec:discussion}

\paragraph{What MoT is and is not.}
MoT is a training-free state-action framework. It is not a correctness classifier, not a dense head-control method, and not a learned policy controller. Its purpose is to map local internal regimes to better default actions with explicit structural safeguards.

\paragraph{Why the design is lightweight.}
The method keeps only five state dimensions and a small action/memory vocabulary. This reduces engineering complexity while preserving core adaptivity for revision-heavy and long-horizon reasoning.

\paragraph{Limitations.}
Performance depends on state coverage quality and calibration of structural budgets. Extremely sparse regime regions can lead to unstable action preference estimates. Preference estimation is observational and may suffer from confounding; MoT uses it as a practical control signal rather than a causal guarantee. The current version is sentence-level and may miss useful token-level cues.

\section{Future Work}
\label{sec:future}

Future directions include: (i) stronger uncertainty calibration for regime assignment, (ii) token-level refinement on top of sentence-level actions, and (iii) optional internal enhancements (including head modulation) as secondary modules rather than the main paradigm carrier.



\bibliographystyle{plain} %%%%% 仅显示数字
% \bibliographystyle{unsrtnat} %%%%% 显示作者信息
\bibliography{reference}




\end{document}